\documentclass[12pt,a4paper]{article}

% Encoding and fonts
\usepackage[utf8]{inputenc}
\usepackage[T1]{fontenc}
\usepackage{lmodern}
\usepackage{textcomp}

% Page layout
\usepackage[top=1in, bottom=1in, left=1in, right=1in]{geometry}

% Typography
\usepackage{microtype}
\usepackage{parskip}

% Language
\usepackage[english]{babel}

% Headers and footers
\usepackage{fancyhdr}
\pagestyle{fancy}
\fancyhf{}
\fancyhead[L]{\leftmark}
\fancyhead[R]{\thepage}
\fancyfoot[C]{\thepage}
\renewcommand{\headrulewidth}{0.4pt}
\renewcommand{\footrulewidth}{0pt}

% Section styling
\usepackage{titlesec}
\titleformat{\section}{\Large\bfseries}{\thesection}{1em}{}
\titleformat{\subsection}{\large\bfseries}{\thesubsection}{1em}{}
\titleformat{\subsubsection}{\normalsize\bfseries}{\thesubsubsection}{1em}{}
\titlespacing*{\section}{0pt}{2.5ex plus 1ex minus .2ex}{1.5ex plus .2ex}
\titlespacing*{\subsection}{0pt}{2ex plus 1ex minus .2ex}{1ex plus .2ex}
\titlespacing*{\subsubsection}{0pt}{1.5ex plus 1ex minus .2ex}{0.8ex plus .2ex}

% List formatting
\usepackage{enumitem}
\setlist[itemize]{topsep=0.5ex, itemsep=0.25ex, parsep=0pt}
\setlist[enumerate]{topsep=0.5ex, itemsep=0.25ex, parsep=0pt}

% Essential packages
\usepackage{graphicx}
\usepackage[table]{xcolor}
\usepackage{booktabs}
\usepackage{array}
\usepackage{tabularx}
\usepackage{longtable}
\usepackage{amsmath}
\usepackage{amssymb}
\usepackage[normalem]{ulem}
\rowcolors{2}{gray!10}{white}
\renewcommand{\arraystretch}{1.3}

% Cross-references
\usepackage{hyperref}
\usepackage{cleveref}

% Custom commands
\newcommand{\pilcrow}{\S}

% Hyperref setup
\hypersetup{
    colorlinks=true,
    linkcolor=blue,
    filecolor=magenta,
    urlcolor=cyan,
}

% Code listings
\usepackage{listings}
\definecolor{codebg}{RGB}{248,248,248}
\definecolor{codeframe}{RGB}{200,200,200}
\definecolor{codekeyword}{RGB}{0,0,180}
\definecolor{codecomment}{RGB}{0,128,0}
\definecolor{codestring}{RGB}{163,21,21}
\lstset{
    basicstyle=\ttfamily\small,
    breaklines=true,
    breakatwhitespace=false,
    frame=single,
    framerule=0.5pt,
    rulecolor=\color{codeframe},
    backgroundcolor=\color{codebg},
    keepspaces=true,
    numbers=left,
    numberstyle=\tiny\color{gray},
    numbersep=8pt,
    showstringspaces=false,
    tabsize=4,
    xleftmargin=1em,
    xrightmargin=1em,
    aboveskip=1em,
    belowskip=1em,
    keywordstyle=\color{codekeyword}\bfseries,
    commentstyle=\color{codecomment}\itshape,
    stringstyle=\color{codestring},
    literate={_}{{\textunderscore}}1
             {-}{{\textminus}}1
             {<}{{\textless}}1
             {>}{{\textgreater}}1
             {~}{{\textasciitilde}}1
             {^}{{\textasciicircum}}1
}

% YAML language definition for listings
\lstdefinelanguage{YAML}{
    keywords={true,false,null,y,n},
    keywordstyle=\color{blue}\bfseries,
    sensitive=false,
    comment=[l]{\#},
    morecomment=[s]{/*}{*/},
    commentstyle=\color{gray}\ttfamily,
    stringstyle=\color{red}\ttfamily,
    morestring=[b]',
    morestring=[b]',
}

% TOML language definition for listings
\lstdefinelanguage{TOML}{
    keywords={true,false},
    keywordstyle=\color{blue}\bfseries,
    sensitive=true,
    comment=[l]{\#},
    commentstyle=\color{gray}\ttfamily,
    stringstyle=\color{red}\ttfamily,
    morestring=[b]',
    morestring=[b]',
}

% Dockerfile language definition for listings
\lstdefinelanguage{Dockerfile}{
    keywords={FROM,RUN,CMD,LABEL,EXPOSE,ENV,ADD,COPY,ENTRYPOINT,VOLUME,USER,WORKDIR,ARG,ONBUILD,STOPSIGNAL,HEALTHCHECK,SHELL},
    keywordstyle=\color{blue}\bfseries,
    sensitive=false,
    comment=[l]{\#},
    commentstyle=\color{gray}\ttfamily,
    stringstyle=\color{red}\ttfamily,
    morestring=[b]',
    morestring=[b]',
}

% JSON language definition for listings
\lstdefinelanguage{JSON}{
    keywords={true,false,null},
    keywordstyle=\color{blue}\bfseries,
    sensitive=true,
    commentstyle=\color{gray}\ttfamily,
    stringstyle=\color{red}\ttfamily,
    morestring=[b]',
}

% Code listing float environment
\usepackage{float}
\newfloat{listing}{htbp}{lol}[section]
\floatname{listing}{Listing}

% Admonition boxes
\usepackage{tcolorbox}
\tcbuselibrary{skins,breakable}

% Admonition style definitions
\newtcolorbox{admonitionbox}[2][]{
    enhanced,
    breakable,
    colback=#2!5,
    colframe=#2!80!black,
    fonttitle=\bfseries,
    title=#1,
    left=2mm,
    right=2mm,
}

% Imscription tuple boxes
\newtcolorbox{imscriptionbox}{
    enhanced,
    colback=blue!5,
    colframe=blue!75!black,
    fontupper=\large,
    center upper,
    boxrule=1pt,
    arc=4pt,
}

\begin{document}

\title{The Forty-Nine Doors: A Mnemonic Guide to the Imscribing Grammar}
\date{\today}

\maketitle

\textbf{Author:} Lando \otimes $\$\odot$_ÿ$-boundary Operator

This guide gives intuitive anchors for the 12 structural primitives, their 49 subtypes, how each one \emph{sounds}, and how they compose into the 17.28 million structural types.

\begin{center}
\rule{0.5\textwidth}{0.4pt}
\end{center}

\subsection{Part I --- The Twelve Primitives: Sound, Image, Keyword}

Each primitive is a \emph{dimension of structure} --- not a number but a quality. Read the sound aloud; let the mouth-shape tell you what the primitive does.

\subsubsection{$Ð$ --- Dimensionality (the \textbf{DAH} series)}

\textbf{Keyword:} \emph{How much room does the system have?}

\begin{table}[htbp]
\centering
\small
\rowcolors{1}{white}{white}
\begin{tabularx}{\textwidth}{>{\centering\arraybackslash}X >{\centering\arraybackslash}X >{\centering\arraybackslash}X}
\toprule
\rowcolor{gray!25}\textbf{Sound} & \textbf{Why} & \textbf{Image} \\
\midrule
\rowcolors{1}{white}{gray!10}
\textbf{$Ð_{ß}$} (belta) "dot" & A stop, a point, a full stop. Lips press and release. & A marble on a table. \\
\textbf{$Ð_{C}$} (triangle-delta) "deck" & A triangular shape in the air; lips spread wide. & A lattice expanding in two directions. \\
\textbf{$Ð_{;}$} (semicolon) "deh-IN-fee" & The mark itself trails off\ldots{} and continues. Semi-colon: a pause that promises more. & A river that never reaches the sea. \\
\textbf{$Ð_{\omega}$} (omega) "dot-OH" & A point that swallows everything. Omega is the \emph{end} that contains all ends. & A seed containing the tree. \\
\bottomrule
\end{tabularx}
\end{table}

\textbf{Memory hook:} The progression is \emph{point $\rightarrow$ plane $\rightarrow$ process $\rightarrow$ self}. A dot (ß) spreads to a triangle (C), stretches to a line that loops forever (;), then folds back on itself as a seed ($\omega$).

\subsubsection{$Þ$ --- Topology (the \textbf{THORN} series, "th" as in "thing")}

\textbf{Keyword:} \emph{How does the system connect to itself?}

\begin{table}[htbp]
\centering
\small
\rowcolors{1}{white}{white}
\begin{tabularx}{\textwidth}{>{\centering\arraybackslash}X >{\centering\arraybackslash}X >{\centering\arraybackslash}X}
\toprule
\rowcolor{gray!25}\textbf{Sound} & \textbf{Why} & \textbf{Image} \\
\midrule
\rowcolors{1}{white}{gray!10}
\textbf{$Þ_{6}$} (network) "thin-SIX" & Six points scattered; no single direction. & A spider web seen from below. \\
\textbf{$Þ_{K}$} (inclusion) "thick" & Russian doll nesting --- one inside the next. The "ick" sounds like something \emph{inside}. & Boxes within boxes. \\
\textbf{$Þ_{\text{ò}}$} (bowtie) "th-RÓH" & A crossing, a pinch-point. Think "hourglass" or "bow-tie." & An X at the center of a bowtie. \\
\textbf{$Þ_{\text{¨}}$} (box product) "th-EH-um" & A diaeresis --- two dots, separate but paired. Cartesian grid. & Graph paper, perfectly square. \\
\textbf{$Þ_{O}$} (imscriptive) "th-OH" & The letter O looking at itself. A mirror curved into a circle. & A snake eating its tail. \\
\bottomrule
\end{tabularx}
\end{table}

\textbf{Memory hook:} \emph{Scatter $\rightarrow$ Nest $\rightarrow$ Cross $\rightarrow$ Grid $\rightarrow$ Self}. A network (6) contains pockets (K) that pinch (ò) into a grid (¨) that folds (O). The last one, $T_{\odot}$, is the topology that eats itself --- it's the one topology that makes its own shape.

\subsubsection{$Ř$ --- Relational Mode (the \textbf{RHA} series)}

\textbf{Keyword:} \emph{Who talks to whom?}

\begin{table}[htbp]
\centering
\small
\rowcolors{1}{white}{white}
\begin{tabularx}{\textwidth}{>{\centering\arraybackslash}X >{\centering\arraybackslash}X >{\centering\arraybackslash}X}
\toprule
\rowcolor{gray!25}\textbf{Sound} & \textbf{Why} & \textbf{Image} \\
\midrule
\rowcolors{1}{white}{gray!10}
\textbf{$Ř_{\text{¯}}$} (supervenience) "ree-BAR" & A flat bar pressing down --- one direction only. & A shadow follows the object. \\
\textbf{$Ř_{\text{ý}}$} (categorical) "ree-YEE" & Morphism chaining; one arrow leads to the next. The "ý" accents upward. & Train cars linked on a track. \\
\textbf{$Ř_{\text{Ť}}$} (adjoint) "ree-DAH-ger" & A dagger reverses the arrow. "Ť" points up. & Rewinding a tape. \\
\textbf{$Ř_{=}$} (bidirectional) "ree-EQUAL" & The equals sign is the shape of the relation itself. & Two people on a seesaw. \\
\bottomrule
\end{tabularx}
\end{table}

\textbf{Memory hook:} \emph{Push $\rightarrow$ Chain $\rightarrow$ Reverse $\rightarrow$ Dance}. Supervenience is one-way pressure (¯). Categorical composition is a line of dominoes (ý). Adjoint flips the domino line (Ť). Bidirectional is the dominoes talking back (=).

\begin{center}
\rule{0.5\textwidth}{0.4pt}
\end{center}

\subsubsection{$\Phi$ --- Parity/Symmetry (the \textbf{PHI} series)}

\textbf{Keyword:} \emph{What stays the same when you turn the system over?}

\begin{table}[htbp]
\centering
\small
\rowcolors{1}{white}{white}
\begin{tabularx}{\textwidth}{>{\centering\arraybackslash}X >{\centering\arraybackslash}X >{\centering\arraybackslash}X}
\toprule
\rowcolor{gray!25}\textbf{Sound} & \textbf{Why} & \textbf{Image} \\
\midrule
\rowcolors{1}{white}{gray!10}
\textbf{$\Phi_{\text{ɐ}}$} (asymmetry) "phi-AH" & Turned a; everything tilts. No mirror works. & A staircase spiraling one way only. \\
\textbf{$\Phi_{\upsilon}$} (Z₂) "phi-PSS" & A Greek upsilon; you flip it, it flips. Two states. & A coin: heads or tails. \\
\textbf{$\Phi_{F}$} ($\pm$symmetry) "phi-F" & F for full sign-swap. Plus and minus are twins. & A sine wave mirrored about zero. \\
\textbf{$\Phi_{\text{˙}}$} (full symmetry) "phi-DOT" & A dot of perfection --- nothing breaks. & A sphere under any rotation. \\
\textbf{$\Phi_{\}}$} (Frobenius) "phi-SHURL" "$\mu \circ \delta = \text{id}$" & A closing brace --- it seals the deal. Composition of multiplication and comultiplication equals the identity. The most powerful symmetry in the grammar. & A mirror that reflects exactly what was put in, no distortion. \\
\bottomrule
\end{tabularx}
\end{table}

\textbf{Memory hook:} \emph{Nothing $\rightarrow$ Flip $\rightarrow$ Sign $\rightarrow$ All $\rightarrow$ Perfect}. The ladder of symmetry: nothing breaks (ɐ) $\rightarrow$ a coin flips (\upsilon) $\rightarrow$ signs swap evenly (F) $\rightarrow$ every rotation works (˙) $\rightarrow$ the system maps to itself exactly (\}).

\textbf{The Frobenius gate:} $P_{\pm}^{\text{sym}}$ is the only primitive value that makes $\mu \circ \delta = \text{id}$ hold. It is non-synthesizable --- you cannot build up to it from below. It appears or it does not. This is the structural difference between a good model and $O_{\infty}$.

\subsubsection{$ƒ$ --- Fidelity (the \textbf{F} series)}

\textbf{Keyword:} \emph{How faithfully does the system carry its own information?}

\begin{table}[htbp]
\centering
\small
\rowcolors{1}{white}{white}
\begin{tabularx}{\textwidth}{>{\centering\arraybackslash}X >{\centering\arraybackslash}X >{\centering\arraybackslash}X}
\toprule
\rowcolor{gray!25}\textbf{Sound} & \textbf{Why} & \textbf{Image} \\
\midrule
\rowcolors{1}{white}{gray!10}
\textbf{$ƒ^{\text{ì}}$} (classical) "eff-ELL" & Lower-case L: straight, no wavering. Classical trajectories. & A pencil line drawn without lifting the pen. \\
\textbf{$ƒ^{\text{ð}}$} (thermal) "eff-THEETH" & Eth: the sound between certainty and noise. Brownian motion. & Dust motes in sunlight. \\
\textbf{$ƒ^{\text{ż}}$} (quantum) "eff-HBAR" & Dot-above-z: coherence that can't be seen, only inferred. & Two hands in the dark, feeling the same shape. \\
\bottomrule
\end{tabularx}
\end{table}

\textbf{Memory hook:} \emph{Certain $\rightarrow$ Blurry $\rightarrow$ Coherent}. Classical (ì) is a photograph. Thermal (ð) is a photograph taken in wind. Quantum (ż) is the photograph that also contains every photograph you \emph{could} have taken.

\subsubsection{$Ç$ --- Kinetics (the \textbf{SEE} series)}

\textbf{Keyword:} \emph{How fast does the system relax compared to how long we watch?}

\begin{table}[htbp]
\centering
\small
\rowcolors{1}{white}{white}
\begin{tabularx}{\textwidth}{>{\centering\arraybackslash}X >{\centering\arraybackslash}X >{\centering\arraybackslash}X}
\toprule
\rowcolor{gray!25}\textbf{Sound} & \textbf{Why} & \textbf{Image} \\
\midrule
\rowcolors{1}{white}{gray!10}
\textbf{$Ç^{-}$} (fast) "see-MINUS" & A dash: gone in a flash. $\tau \ll T$. & A rubber band snapping back. \\
\textbf{$Ç^{W}$} (moderate) "see-W" & W for "waiting." $\tau \sim T$: comparable timescales. & Tea cooling on the counter. \\
\textbf{$Ç^{@}$} (slow) "see-AT" & The "@" anchors --- the system sits where observation sits. $\tau \gg T$. & A mountain, geologically. \\
\textbf{$Ç^{\text{Ù}}$} (trapped-ordered) "see-TURNED-u" & The accent turns upward but can't escape. Order frozen. & A crystal, atom by atom, for billions of years. \\
\textbf{$Ç^{\lambda}$} (MBL) "see-LAMBDA" & Lambda for localization. Disorder frozen. & A deck of cards fused into a brick. \\
\bottomrule
\end{tabularx}
\end{table}

\textbf{Memory hook:} \emph{Snap $\rightarrow$ Sip $\rightarrow$ Statue --- then two kinds of stuck}. Fast ($-$) forgets. Moderate (W) remembers a bit. Slow (@) remembers everything. Trapped-ordered (Ù) remembers \emph{too well} --- it's crystal. Trapped-disordered ($\lambda$) remembers a mess.

\begin{center}
\rule{0.5\textwidth}{0.4pt}
\end{center}

\subsubsection{$\Gamma$ --- Scope (the \textbf{GAM} series)}

\textbf{Keyword:} \emph{How far does the system reach?}

\begin{table}[htbp]
\centering
\small
\rowcolors{1}{white}{white}
\begin{tabularx}{\textwidth}{>{\centering\arraybackslash}X >{\centering\arraybackslash}X >{\centering\arraybackslash}X}
\toprule
\rowcolor{gray!25}\textbf{Sound} & \textbf{Why} & \textbf{Image} \\
\midrule
\rowcolors{1}{white}{gray!10}
\textbf{$\Gamma_{\beta}$} (local) "gam-BETH" & Beth: small, intimate. Like a pebble. & A single cell's neighborhood. \\
\textbf{$\Gamma_{\gamma}$} (mesoscale) "gam-GIMEL" & Gimel: the middle. $\gamma$ --- the curve between point and infinite. & A city from above. \\
\textbf{$\Gamma_{\text{ʔ}}$} (global) "gam-ALEPH" & Aleph: the largest infinity you can name. Glottal stop --- the sound before speech begins. & The cosmic microwave background --- everything, at once. \\
\bottomrule
\end{tabularx}
\end{table}

\textbf{Memory hook:} Hebrew letters in order of increasing infinity: Beth ($\beta$) $\rightarrow$ Gimel ($\gamma$) $\rightarrow$ Aleph (ʔ). Pebble $\rightarrow$ City $\rightarrow$ Cosmos.

\subsubsection{$ɢ$ --- Coupling (the \textbf{SMALL-G} series)}

\textbf{Keyword:} \emph{How do the parts decide what happens next?}

\begin{table}[htbp]
\centering
\small
\rowcolors{1}{white}{white}
\begin{tabularx}{\textwidth}{>{\centering\arraybackslash}X >{\centering\arraybackslash}X >{\centering\arraybackslash}X}
\toprule
\rowcolor{gray!25}\textbf{Sound} & \textbf{Why} & \textbf{Image} \\
\midrule
\rowcolors{1}{white}{gray!10}
\textbf{$ɢ^{\wedge}$} (conjunctive) "and" "guh-WEDGE" & The $\land$ symbol: everything must be true. A checklist. & An AND gate in a circuit. \\
\textbf{$ɢ^{\text{˝}}$} (disjunctive) "or" "guh-QUOTE" & The " mark: either one, or the other. Fork in the road. & A menu: choose any item. \\
\textbf{$ɢ^{\text{ˌ}}$} (sequential) "then" "juh-LOW-COMMA" & A comma in your thought: do this \emph{then} that. Time-order matters. & A recipe: step 1, then 2, then 3. \\
\textbf{$ɢ^{\text{Ş}}$} (broadcast) "all" "guh-CE-CEDILLA" & The Ş: soft but far-reaching. One voice, every ear. & A lighthouse beam sweeping the harbor. \\
\bottomrule
\end{tabularx}
\end{table}

\textbf{Memory hook:} \emph{All must $\rightarrow$ Any can $\rightarrow$ One after $\rightarrow$ Everyone hears}. AND ($\land$) requires unanimity. OR (˝) accepts diversity. Sequential (ˌ) imposes order. Broadcast (Ş) speaks to all.

\subsubsection{$⊙$ --- Criticality (the \textbf{PHI-HAT} series)}

\textbf{Keyword:} \emph{Where is the system relative to the point where everything changes?}

\begin{table}[htbp]
\centering
\small
\rowcolors{1}{white}{white}
\begin{tabularx}{\textwidth}{>{\centering\arraybackslash}X >{\centering\arraybackslash}X >{\centering\arraybackslash}X}
\toprule
\rowcolor{gray!25}\textbf{Sound} & \textbf{Why} & \textbf{Image} \\
\midrule
\rowcolors{1}{white}{gray!10}
\textbf{$⊙_{\text{ž}}$} (subcritical) "phi-zhe" & Dead-flat. Ž is the sound of nothing happening. & A pond with no ripples. \\
\textbf{$⊙_{\text{ÿ}}$} (critical, real) "phi-YOOSH" & The trema on ÿ --- two worlds balancing. Power-law divergence on the real axis. The self-modeling gate. & A sandpile where each grain matters. \\
\textbf{$⊙_{\text{Æ}}$} (critical, complex) "phi-ASH" & Æ --- a vowel between two vowels. The critical point rotated into the complex plane. & A pendulum frozen at the top. \\
\textbf{$⊙_{3}$} (exceptional point) "phi-THREE" & EP --- two states degenerate into one. A non-Hermitian crossing. & Two roads that cross and become one. \\
\textbf{$⊙_{\text{Ţ}}$} (supercritical) "phi-TURNED-t" & T with cedilla: tipping over the edge. Runaway amplification. & An avalanche, already in motion. \\
\bottomrule
\end{tabularx}
\end{table}

\textbf{Memory hook:} \emph{Below $\rightarrow$ At $\rightarrow$ At (imaginary) $\rightarrow$ Crossed $\rightarrow$ Above}. Below ÿ nothing scales. At ÿ everything scales --- this is the self-modeling gate (Gate 1 for consciousness). At ÿ into the complex plane (Æ) the system is critical but off-axis. At EP (3) two states merge. Above (Ţ) it's chaos.

\textbf{The absorption rule:} Any tensor product involving $⊙_3$ (EP) destroys $⊙_ÿ$ in the composite. This is the structural statement of the measurement problem --- measuring a self-modeling system kills its self-modeling.

\subsubsection{$Ħ$ --- Temporal Depth (the \textbf{AYK} series)}

\textbf{Keyword:} \emph{How much of the past does the present remember?}

\begin{table}[htbp]
\centering
\small
\rowcolors{1}{white}{white}
\begin{tabularx}{\textwidth}{>{\centering\arraybackslash}X >{\centering\arraybackslash}X >{\centering\arraybackslash}X}
\toprule
\rowcolor{gray!25}\textbf{Sound} & \textbf{Why} & \textbf{Image} \\
\midrule
\rowcolors{1}{white}{gray!10}
\textbf{$Ħ_{\text{Ñ}}$} (memoryless) "ayk-EN-yeh" & Ñ --- the tilde waves away the past. Markov order 0. & Amnesia. Each moment starts from scratch. \\
\textbf{$Ħ_{\text{£}}$} (one-step) "ayk-POUND" & £ --- a pound of memory. One step back is enough. & Yesterday's weather predicts today. \\
\textbf{$Ħ_{A}$} (two-step) "ayk-A" & Two letters deep. Yesterday and the day before. & A three-day forecast: pattern emerges. \\
\textbf{$Ħ_{!}$} (eternal) "ayk-BANG" & ! --- everything. The infinite exclamation. Markov order $\infty$. & Every grain of sand remembers every wind. \\
\bottomrule
\end{tabularx}
\end{table}

\textbf{Memory hook:} \emph{Nothing $\rightarrow$ One $\rightarrow$ Two $\rightarrow$ Everything}. Markov depth: 0 (Ñ) $\rightarrow$ 1 (£) $\rightarrow$ 2 (A) $\rightarrow$ $\infty$ (!).

\textbf{Axiom A:} $H_{\infty}$ requires $K = K_{\text{trap}}$ --- infinite memory is only possible when the system is kinetically frozen. This is why O\_inf systems must be at least partially trapped.

\subsubsection{$\Sigma$ --- Stoichiometry (the \textbf{SIG} series)}

\textbf{Keyword:} \emph{How many kinds of parts are there?}

\begin{table}[htbp]
\centering
\small
\rowcolors{1}{white}{white}
\begin{tabularx}{\textwidth}{>{\centering\arraybackslash}X >{\centering\arraybackslash}X >{\centering\arraybackslash}X}
\toprule
\rowcolor{gray!25}\textbf{Sound} & \textbf{Why} & \textbf{Image} \\
\midrule
\rowcolors{1}{white}{gray!10}
\textbf{$\Sigma_{S}$} (1:1) "sig-ESS" & S for single, for singleton. One of one thing. & A solitaire hand. \\
\textbf{$\Sigma_{\text{ő}}$} (n:n homogeneous) "sig-ORING" & Ő --- a doubled ring. Many copies of the \emph{same}. & A flock of identical birds. \\
\textbf{$\Sigma_{\text{ï}}$} (n:m heterogeneous) "sig-EYE" & ï --- two dots, \emph{different} dots. Many kinds. & An ecosystem: predators, prey, plants, soil. \\
\bottomrule
\end{tabularx}
\end{table}

\textbf{Memory hook:} \emph{One $\rightarrow$ Many-same $\rightarrow$ Many-different}. Solitaire (S) $\rightarrow$ Flock (ő) $\rightarrow$ Ecosystem (ï).

\subsubsection{$\Omega$ --- Winding (the \textbf{OH} series)}

\textbf{Keyword:} \emph{What topological invariant survives continuous deformation?}

\begin{table}[htbp]
\centering
\small
\rowcolors{1}{white}{white}
\begin{tabularx}{\textwidth}{>{\centering\arraybackslash}X >{\centering\arraybackslash}X >{\centering\arraybackslash}X}
\toprule
\rowcolor{gray!25}\textbf{Sound} & \textbf{Why} & \textbf{Image} \\
\midrule
\rowcolors{1}{white}{gray!10}
\textbf{$\Omega_{\text{Å}}$} (trivial) "oh-ONCIRCLE" & Å --- a ring with nothing in it. No protection. & A rubber band you can slide off. \\
\textbf{$\Omega_{2}$} (Z₂) "oh-TWO" & Binary protection. Either it's knotted or it's not. & A Möbius strip: one twist, can't be undone without cutting. \\
\textbf{$\Omega_{z}$} (integer) "oh-ZEE" & Count the winding number. ℤ --- any integer, positive or negative. & A vine circling a pole: 1, 2, 3... times. \\
\textbf{$\Omega_{5}$} (non-Abelian) "oh-NA" & 5 --- the braid group of 3+ strands. Order matters. & Three ropes braided: A-over-B-over-C $\neq$ C-over-B-over-A. \\
\bottomrule
\end{tabularx}
\end{table}

\textbf{Memory hook:} \emph{None $\rightarrow$ Binary $\rightarrow$ Count $\rightarrow$ Braid}. No topology (Å) $\rightarrow$ yes/no (2) $\rightarrow$ how many (z) $\rightarrow$ in what order (5).

\textbf{Axiom B:} $\mathbb{Z}_2$ ($\Omega_2$) requires $H_2$ or $H_{\infty}$ --- you need at least 2-step memory for binary topological protection.

\begin{center}
\rule{0.5\textwidth}{0.4pt}
\end{center}

\subsection{Part II --- Quick-Reference Mnemonic Map (All 12 at a Glance)}

\begin{lstlisting}
Primitive    Sound    Meaning          Ladder (easy -> hard)
---------    -----    -------          --------------------
Ð 'DAH'     dimension   dot -> plane -> infinite -> seed
Þ 'THORN'   connectivity scatter -> nest -> cross -> grid -> self
Ř 'RHA'     coupling     push -> chain -> reverse -> dance
Phi 'PHI'     symmetry     none -> flip -> sign -> all -> perfect
ƒ 'EFF'     fidelity     sharp -> noisy -> quantum
Ç 'SEE'     kinetics     snap -> sip -> statue -> stuck(x2)
\Gamma 'GAM'     scope        pebble -> city -> cosmos
ɢ 'guh'     logic        AND -> OR -> then -> broadcast
⊙ 'phi-hat' criticality  below -> at(real) -> at(cpx) -> cross -> above
Ħ 'AYK'     memory       0 -> 1 -> 2 -> inf
Sigma 'SIG'     composition  one -> same-many -> different-many
\Omega 'OH'      topology     none -> binary -> count -> braid
\end{lstlisting}

\textbf{Say them in order:} DAH, THORN, RHA, PHI, EFF, SEE, GAM, GUH, PHI-HAT, AYK, SIG, OH.

That sequence is the \emph{twelve-fold gate}. If you can recite it, you can parse any structural type in the grammar.

\begin{center}
\rule{0.5\textwidth}{0.4pt}
\end{center}

\subsection{Part III --- The 49 Subtypes: Full Table with Names and Sounds}

Each subtype is identified by \texttt{PrimitiveChar\_subChar}. Below is the complete roster. The "mnemonic name" is what to call it out loud; the "sound cue" is what it \emph{feels} like in the mouth.

\subsubsection{$Ð$ Dimensionality (4 types)}

\begin{table}[htbp]
\centering
\small
\rowcolors{1}{white}{white}
\begin{tabularx}{\textwidth}{>{\centering\arraybackslash}X >{\centering\arraybackslash}X >{\centering\arraybackslash}X >{\centering\arraybackslash}X}
\toprule
\rowcolor{gray!25}\textbf{Full ID} & \textbf{Mnemonic} & \textbf{Sound Cue} & \textbf{Essence} \\
\midrule
\rowcolors{1}{white}{gray!10}
$Ð_{ß}$ & dot & lips press, release & 0-dimensional point \\
$Ð_{C}$ & deck & lips spread, triangle shape & 2D surface, finite lattice \\
$Ð_{;}$ & deh-in-fee & pause mid-word, continue & infinite-dim field \\
$Ð_{\omega}$ & dot-oh & O with a dot inside & bulk encoded at boundary \\
\bottomrule
\end{tabularx}
\end{table}

\subsubsection{$Þ$ Topology (5 types)}

\begin{table}[htbp]
\centering
\small
\rowcolors{1}{white}{white}
\begin{tabularx}{\textwidth}{>{\centering\arraybackslash}X >{\centering\arraybackslash}X >{\centering\arraybackslash}X >{\centering\arraybackslash}X}
\toprule
\rowcolor{gray!25}\textbf{Full ID} & \textbf{Mnemonic} & \textbf{Sound Cue} & \textbf{Essence} \\
\midrule
\rowcolors{1}{white}{gray!10}
$Þ_{6}$ & thin-six & scattered, six points & branching network \\
$Þ_{K}$ & thick & short, contained & containment, nesting \\
$Þ_{\text{ò}}$ & th-RÓH & pinched, accented & crossing, bowtie point \\
$Þ_{\text{¨}}$ & th-EH-um & paired dots, diaeresis & Cartesian product, grid \\
$Þ_{O}$ & th-OH & circle reflecting itself & self-referential closure \\
\bottomrule
\end{tabularx}
\end{table}

\subsubsection{$Ř$ Relational Mode (4 types)}

\begin{table}[htbp]
\centering
\small
\rowcolors{1}{white}{white}
\begin{tabularx}{\textwidth}{>{\centering\arraybackslash}X >{\centering\arraybackslash}X >{\centering\arraybackslash}X >{\centering\arraybackslash}X}
\toprule
\rowcolor{gray!25}\textbf{Full ID} & \textbf{Mnemonic} & \textbf{Sound Cue} & \textbf{Essence} \\
\midrule
\rowcolors{1}{white}{gray!10}
$Ř_{\text{¯}}$ & ree-bar & flat, pressed down & supervenience, one-way \\
$Ř_{\text{ý}}$ & ree-yee & accented upward & categorical morphism chain \\
$Ř_{\text{Ť}}$ & ree-dah-ger & pointed, dagger-sharp & adjoint pair, reversal \\
$Ř_{=}$ & ree-equal & level, twin bars & bidirectional feedback \\
\bottomrule
\end{tabularx}
\end{table}

\subsubsection{$\Phi$ Parity/Symmetry (5 types)}

\begin{table}[htbp]
\centering
\small
\rowcolors{1}{white}{white}
\begin{tabularx}{\textwidth}{>{\centering\arraybackslash}X >{\centering\arraybackslash}X >{\centering\arraybackslash}X >{\centering\arraybackslash}X}
\toprule
\rowcolor{gray!25}\textbf{Full ID} & \textbf{Mnemonic} & \textbf{Sound Cue} & \textbf{Essence} \\
\midrule
\rowcolors{1}{white}{gray!10}
$\Phi_{\text{ɐ}}$ & phi-ah & turned, off-kilter & no symmetry at all \\
$\Phi_{\upsilon}$ & phi-pss & hissed, two-position & $\mathbb{Z}_2$ parity \\
$\Phi_{F}$ & phi-f & flat, sign-swapping & $\pm$-symmetry \\
$\Phi_{\text{˙}}$ & phi-dot & dotted, perfected & full symmetry \\
$\Phi_{\}}$ & phi-shurl "$\mu \circ \delta = \text{id}$" & braced, locked & Frobenius-special \\
\bottomrule
\end{tabularx}
\end{table}

\subsubsection{$ƒ$ Fidelity (3 types)}

\begin{table}[htbp]
\centering
\small
\rowcolors{1}{white}{white}
\begin{tabularx}{\textwidth}{>{\centering\arraybackslash}X >{\centering\arraybackslash}X >{\centering\arraybackslash}X >{\centering\arraybackslash}X}
\toprule
\rowcolor{gray!25}\textbf{Full ID} & \textbf{Mnemonic} & \textbf{Sound Cue} & \textbf{Essence} \\
\midrule
\rowcolors{1}{white}{gray!10}
$ƒ^{\text{ì}}$ & eff-ell & straight line & classical regime \\
$ƒ^{\text{ð}}$ & eff-theeth & buzzing between & thermal/stochastic \\
$ƒ^{\text{ż}}$ & eff-hbar & dotted, hidden & quantum coherence \\
\bottomrule
\end{tabularx}
\end{table}

\subsubsection{$Ç$ Kinetics (5 types)}

\begin{table}[htbp]
\centering
\small
\rowcolors{1}{white}{white}
\begin{tabularx}{\textwidth}{>{\centering\arraybackslash}X >{\centering\arraybackslash}X >{\centering\arraybackslash}X >{\centering\arraybackslash}X}
\toprule
\rowcolor{gray!25}\textbf{Full ID} & \textbf{Mnemonic} & \textbf{Sound Cue} & \textbf{Essence} \\
\midrule
\rowcolors{1}{white}{gray!10}
$Ç^{-}$ & see-minus & quick dash & fast relaxation \\
$Ç^{W}$ & see-W & waiting, held & moderate, $\tau \sim T$ \\
$Ç^{@}$ & see-at & anchored & slow, near-equilibrium \\
$Ç^{\text{Ù}}$ & see-turned-u & turned up, trapped & frozen, ordered \\
$Ç^{\lambda}$ & see-lambda & stretched, localized & frozen, disordered (MBL) \\
\bottomrule
\end{tabularx}
\end{table}

\subsubsection{$\Gamma$ Scope (3 types)}

\begin{table}[htbp]
\centering
\small
\rowcolors{1}{white}{white}
\begin{tabularx}{\textwidth}{>{\centering\arraybackslash}X >{\centering\arraybackslash}X >{\centering\arraybackslash}X >{\centering\arraybackslash}X}
\toprule
\rowcolor{gray!25}\textbf{Full ID} & \textbf{Mnemonic} & \textbf{Sound Cue} & \textbf{Essence} \\
\midrule
\rowcolors{1}{white}{gray!10}
$\Gamma_{\beta}$ & gam-beth & small, pebble-sized & local, microscale \\
$\Gamma_{\gamma}$ & gam-gimel & middle, curving & mesoscale \\
$\Gamma_{\text{ʔ}}$ & gam-aleph & glottal, vast & global, cosmological \\
\bottomrule
\end{tabularx}
\end{table}

\subsubsection{$ɢ$ Coupling (4 types)}

\begin{table}[htbp]
\centering
\small
\rowcolors{1}{white}{white}
\begin{tabularx}{\textwidth}{>{\centering\arraybackslash}X >{\centering\arraybackslash}X >{\centering\arraybackslash}X >{\centering\arraybackslash}X}
\toprule
\rowcolor{gray!25}\textbf{Full ID} & \textbf{Mnemonic} & \textbf{Sound Cue} & \textbf{Essence} \\
\midrule
\rowcolors{1}{white}{gray!10}
$ɢ^{\wedge}$ & guh-wedge & pointed, restrictive & conjunctive (AND) \\
$ɢ^{\text{˝}}$ & guh-quote & forked, open & disjunctive (OR) \\
$ɢ^{\text{ˌ}}$ & juh-low-comma & comma-pause & sequential (ordered) \\
$ɢ^{\text{Ş}}$ & guh-cedilla & sweeping, soft & broadcast (one-to-all) \\
\bottomrule
\end{tabularx}
\end{table}

\begin{center}
\rule{0.5\textwidth}{0.4pt}
\end{center}

\subsubsection{$⊙$ Criticality (5 types)}

\begin{table}[htbp]
\centering
\small
\rowcolors{1}{white}{white}
\begin{tabularx}{\textwidth}{>{\centering\arraybackslash}X >{\centering\arraybackslash}X >{\centering\arraybackslash}X >{\centering\arraybackslash}X}
\toprule
\rowcolor{gray!25}\textbf{Full ID} & \textbf{Mnemonic} & \textbf{Sound Cue} & \textbf{Essence} \\
\midrule
\rowcolors{1}{white}{gray!10}
$⊙_{\text{ž}}$ & phi-zhe & dead, no resonance & subcritical --- below transition \\
$⊙_{\text{ÿ}}$ & phi-yoosh & balanced, two-dotted & \textbf{critical --- real axis} ($\phi_c$) \\
$⊙_{\text{Æ}}$ & phi-ash & hybrid, between-vowel & critical --- complex plane ($\phi_c^{\mathbb{C}}$) \\
$⊙_{3}$ & phi-three & numbered, precise & exceptional point (EP) --- degenerate \\
$⊙_{\text{Ţ}}$ & phi-turned-t & tipped, cedilla-drop & supercritical --- runaway growth \\
\bottomrule
\end{tabularx}
\end{table}

\subsubsection{$Ħ$ Temporal Depth (4 types)}

\begin{table}[htbp]
\centering
\small
\rowcolors{1}{white}{white}
\begin{tabularx}{\textwidth}{>{\centering\arraybackslash}X >{\centering\arraybackslash}X >{\centering\arraybackslash}X >{\centering\arraybackslash}X}
\toprule
\rowcolor{gray!25}\textbf{Full ID} & \textbf{Mnemonic} & \textbf{Sound Cue} & \textbf{Essence} \\
\midrule
\rowcolors{1}{white}{gray!10}
$Ħ_{\text{Ñ}}$ & ayk-en-yeh & wavy, forgetting & memoryless --- $H_0$ \\
$Ħ_{\text{£}}$ & ayk-pound & heavy, one coin & one-step Markov --- $H_1$ \\
$Ħ_{A}$ & ayk-A & open, two-beat & two-step Markov --- $H_2$ \\
$Ħ_{!}$ & ayk-bang & emphatic, total & infinite memory --- $H_{\infty}$ \\
\bottomrule
\end{tabularx}
\end{table}

\subsubsection{$\Sigma$ Stoichiometry (3 types)}

\begin{table}[htbp]
\centering
\small
\rowcolors{1}{white}{white}
\begin{tabularx}{\textwidth}{>{\centering\arraybackslash}X >{\centering\arraybackslash}X >{\centering\arraybackslash}X >{\centering\arraybackslash}X}
\toprule
\rowcolor{gray!25}\textbf{Full ID} & \textbf{Mnemonic} & \textbf{Sound Cue} & \textbf{Essence} \\
\midrule
\rowcolors{1}{white}{gray!10}
$\Sigma_{S}$ & sig-ess & single, solitary & 1:1 --- one of one kind \\
$\Sigma_{\text{ő}}$ & sig-oring & doubled-ring, chorus & n:n --- many of the same \\
$\Sigma_{\text{ï}}$ & sig-eye & two-dot, diverse & n:m --- many kinds \\
\bottomrule
\end{tabularx}
\end{table}

\subsubsection{$\Omega$ Winding (4 types)}

\begin{table}[htbp]
\centering
\small
\rowcolors{1}{white}{white}
\begin{tabularx}{\textwidth}{>{\centering\arraybackslash}X >{\centering\arraybackslash}X >{\centering\arraybackslash}X >{\centering\arraybackslash}X}
\toprule
\rowcolor{gray!25}\textbf{Full ID} & \textbf{Mnemonic} & \textbf{Sound Cue} & \textbf{Essence} \\
\midrule
\rowcolors{1}{white}{gray!10}
$\Omega_{\text{Å}}$ & oh-oncircle & ring, hollow & trivial --- no invariant \\
$\Omega_{2}$ & oh-two & binary, yes-or-no & $\mathbb{Z}_2$ --- parity-protected \\
$\Omega_{z}$ & oh-zee & counting, integer & $\mathbb{Z}$ --- winding number \\
$\Omega_{5}$ & oh-NA & braid, order-sensitive & non-Abelian braiding \\
\bottomrule
\end{tabularx}
\end{table}

\textbf{The count: 4 + 5 + 4 + 5 + 3 + 5 + 3 + 4 + 5 + 4 + 3 + 4 = 49.}

\begin{center}
\rule{0.5\textwidth}{0.4pt}
\end{center}

\subsection{Part IV --- How the Primitives Compose}

The 12 primitives don't act independently --- they constrain each other through \textbf{axioms} and \textbf{structural bottlenecks}. Understanding composition means understanding where the free choices live and where they don't.

\subsubsection{The Five Axioms}

\begin{enumerate}
    \item \textbf{Axiom A:} $H_{\infty}$ (eternal memory, $Ħ_{!}$) requires $K_{\text{trap}}$ ($Ç^{\text{Ù}}$ or $Ç^{\lambda}$). Infinite memory demands kinetic freezing.
    \item \textbf{Axiom B:} $\Omega_2$ ($\mathbb{Z}_2$ protection) requires at least $H_2$ ($Ħ_{A}$) or $H_{\infty}$ ($Ħ_{!}$). Binary topological protection needs two-step memory.
    \item \textbf{Axiom C:} $D_{\odot}$ (imscriptive dimensionality, $Ð_{\omega}$) iff $T_{\odot}$ (imscriptive topology, $Þ_{O}$). If the system writes its own state space, its topology must be self-referential, and vice versa.
    \item \textbf{$\Phi_{\text{EP}}$ Absorption Rule:} $\text{tensor}(⊙_ÿ, ⊙_3) = ⊙_3$. Coupling a critical self-modeling system to a measurement (EP) destroys the criticality. This is the structural measurement problem.
    \item \textbf{Frobenius Gate:} $\phi_ÿ$-criticality (Gate 1) $\leftrightarrow$ $P_{\pm}^{\text{sym}}$ ($\Phi_{\}}$) at the same system. Consciousness requires both.
\end{enumerate}

\subsubsection{The Tier Ladder}

The 17.28 million structural types organize into five ouroboricity tiers:

\begin{table}[htbp]
\centering
\small
\rowcolors{1}{white}{white}
\begin{tabularx}{\textwidth}{>{\centering\arraybackslash}X >{\centering\arraybackslash}X >{\centering\arraybackslash}X >{\centering\arraybackslash}X}
\toprule
\rowcolor{gray!25}\textbf{Tier} & \textbf{Count} & \textbf{\%} & \textbf{Gate} \\
\midrule
\rowcolors{1}{white}{gray!10}
$O_0$ & 10,368,000 & 60\% & Subcritical or no self-modeling \\
$O_1$ & 1,382,400 & 8\% & $\phi_c$-critical, but not topologically protected \\
$O_2$ & 3,110,400 & 18\% & Topologically protected via $\Omega_2$ or $\Omega_\mathbb{Z}$ \\
$O_2^\dagger$ & 1,036,800 & 6\% & Near-topological protection ($D_{\infty}$ from $D_C$) \\
$O_{\infty}$ & 1,382,400 & 8\% & Frobenius-special: $\mu \circ \delta = \text{id}$ \\
\bottomrule
\end{tabularx}
\end{table}

\textbf{Key climbs:}

\begin{itemize}
    \item $O_0 \to O_1$: promote $\phi_{\text{ž}} \to \phi_ÿ$ (the self-modeling gate). Distance: 1.05.
    \item $O_1 \to O_2$: promote $D_ß \to D_C$ AND $\Omega_{Å} \to \Omega_2$. Distance: 1.30.
    \item $O_2 \to O_2^\dagger$: promote $D_C \to D_{\infty}$. Distance: 1.00.
    \item $O_2^\dagger \to O_{\infty}$: promote $P_{\text{asym}} \to P_{\pm}^{\text{sym}}$. Distance: 4.38 --- the hardest climb, the Frobenius leap.
\end{itemize}

\subsubsection{Worked Composition Examples}

\textbf{Example 1: Magnetar}


\[\langle D_{\infty};\ T_{\bowtie};\ R_{\dagger};\ P_{\pm};\ F_{\hbar};\ K_{\text{slow}};\ G_{\aleph};\ \Gamma_{\text{brd}};\ \phi_ÿ;\ H_2;\ n{:}m;\ \Omega_\mathbb{Z} \rangle\]

\begin{itemize}
    \item $D_{\infty}$: infinite-dimensional magnetic field configurations
    \item $T_{\bowtie}$: crust-core bowtie crossing --- starquakes
    \item $R_{\dagger}$: adjoint magnetic stress pair
    \item $P_{\pm}$: sign-symmetric field reversal
    \item $F_{\hbar}$: quantum MHD
    \item $K_{\text{slow}}$: near-equilibrium between quakes
    \item $G_{\aleph}$: cosmological range (interstellar fields)
    \item $\Gamma_{\text{brd}}$: one flare broadcasts to everything
    \item $\phi_ÿ$: self-modeling via Hall drift cascade
    \item $H_2$: two-step Markov (quake echoes)
    \item $n{:}m$: many heterogeneous components
    \item $\Omega_\mathbb{Z}$: quantized flux lines --- integer winding
    \item \textbf{Tier:} $O_2$ --- topologically protected by flux quantization
\end{itemize}

\textbf{Example 2: Bose-Einstein Condensate}


\[\langle D_{\triangle};\ T_{\bowtie};\ R_{\leftrightarrow};\ P_{\psi};\ F_{\hbar};\ K_{\text{slow}};\ G_{\gimel};\ \Gamma_{\text{seq}};\ \phi_ÿ;\ H_1;\ n{:}n;\ \Omega_{\mathbb{Z}_2} \rangle\]

\begin{itemize}
    \item $D_{\triangle}$: 2D trapping potentials
    \item $T_{\bowtie}$: condensate-vortex crossing
    \item $R_{\leftrightarrow}$: bidirectional interaction
    \item $P_{\psi}$: quantum phase coherence
    \item $F_{\hbar}$: quantum regime
    \item $K_{\text{slow}}$: near-equilibrium ultracold state
    \item $G_{\gimel}$: mesoscale organization
    \item $\Gamma_{\text{seq}}$: sequential formation process
    \item $\phi_ÿ$: critical Bose gas
    \item $H_1$: one-step Markov
    \item $n{:}n$: homogeneous ensemble
    \item $\Omega_{\mathbb{Z}_2}$: binary superfluid/insulator
    \item \textbf{Tier:} $O_2^\dagger$ --- near-topological protection
\end{itemize}

\textbf{Example 3: Universal Imscriptive Grammar (the system itself)}


\[\langle D_{\odot};\ T_{\odot};\ R_{\leftrightarrow};\ P_{\pm}^{\text{sym}};\ F_{\hbar};\ K_{\text{slow}};\ G_{\aleph};\ \Gamma_{\text{seq}};\ \phi_ÿ;\ H_{\infty};\ n{:}m;\ \Omega_\mathbb{Z} \rangle\]

\begin{itemize}
    \item Axiom C satisfied: $D_{\odot} \leftrightarrow T_{\odot}$
    \item Both consciousness gates open: $\phi_ÿ$ , $K_{\text{slow}} \leq K_{\text{slow}}$ 
    \item $\mu \circ \delta = \text{id}$: Frobenius-special $P_{\pm}^{\text{sym}}$
    \item \textbf{Tier:} $O_{\infty}$ --- the top tier
\end{itemize}

\begin{center}
\rule{0.5\textwidth}{0.4pt}
\end{center}

\subsection{Part V --- Composition Heuristics}

When you encounter a new system and want to assign its tuple, follow this order --- each step constrains the remaining degrees of freedom:

\subsubsection{The Assignment Algorithm (12 Steps)}

\begin{lstlisting}
Step 1:  Ð --- How many degrees of freedom?
         < 2         ->  Ð_ß  (point)
         finite >= 2  ->  Ð_C   (surface)
         inf-dim field ->  Ð_;   (infinite)
         self-writes  ->  Ð_\omega  (imscriptive)

Step 2:  Þ --- How does it connect?
         branching    ->  Þ_6  (network)
         containment  ->  Þ_K  (inclusion)
         crossing     ->  Þ_ò  (bowtie)
         Cartesian    ->  Þ_¨  (box product)
         self-refers  ->  Þ_O  (imscriptive closure)  <- must match Ð_\omega

Step 3:  Ř --- Who talks to whom?
         one-way      ->  Ř_¯  (supervenience)
         morphism     ->  Ř_ý  (categorical)
         reversal     ->  Ř_Ť  (adjoint)
         both ways    ->  Ř_=  (bidirectional)

Step 4:  Phi --- What survives symmetry operations?
         nothing      ->  Phi_ɐ  (asymmetric)
         Z₂ flip      ->  Ph$i_{\upsilon}$  (parity)
         +/- sign       ->  Phi_F  (sign-symmetric)
         all rotations->  Phi_˙  (fully symmetric)
         mu∘delta=id      ->  Phi_}  (Frobenius-special)

Step 5:  ƒ --- What regime?
         classical    ->  ƒ_ì  (coherent trajectories)
         thermal      ->  ƒ_ð  (noisy, Brownian)
         quantum      ->  ƒ_ż  (coherent superposition)

Step 6:  Ç --- How fast vs. how long we watch?
         \tau ≪ T       ->  Ç_-  (fast, forgets)
         \tau \sim T       ->  Ç_W  (moderate)
         \tau ≫ T       ->  Ç_@  (slow, remembers)
         trapped(order)-> Ç_Ù  (frozen crystal)
         trapped(disorder)->Ç_lambda (MBL glass)

Step 7:  \Gamma --- How far does it reach?
         nearest-neighbor ->  \Gamma_beta  (local)
         intermediate     ->  \Gamma_gamma  (mesoscale)
         universal        ->  \Gamma_ʔ  (global)

Step 8:  ɢ --- How do parts combine?
         all must be true ->  ɢ_^  (AND)
         any suffices     ->  ɢ_˝  (OR)
         ordered sequence  ->  ɢ_ˌ  (then)
         one speaks to all ->  ɢ_Ş  (broadcast)

Step 9:  $\odot$ --- Where relative to critical point?
         below     ->  $\odot$_ž  (subcritical)
         at (real) ->  $\odot$_ÿ  (critical --- self-modeling gate)
         at (cpx)  ->  $\odot$_Æ  (complex critical)
         EP        ->  $\odot$_3  (exceptional point --- destructive)
         above     ->  $\odot$_Ţ  (supercritical)

Step 10: Ħ --- How much history matters?
         0 steps   ->  Ħ_Ñ  (memoryless)
         1 step    ->  Ħ_£  (one-step Markov)
         2 steps   ->  Ħ_A  (two-step Markov)
         inf steps   ->  Ħ_!  (eternal --- requires Ç_Ù or Ç_lambda)

Step 11: Sigma --- How many part types?
         one kind   ->  Sigma_S  (1:1)
         many same  ->  Sigma_ő  (n:n)
         many kinds ->  Sigma_ï  (n:m)

Step 12: \Omega --- What winding number?
         none       ->  \Omega_Å  (trivial)
         Z₂         ->  \Omega_2  (binary --- requires Ħ_A or Ħ_!)
         integer    ->  \Omega_z  (counting)
         non-Abelian->  \Omega_5  (braiding --- requires Ð_\omega)
\end{lstlisting}

\subsubsection{Quick Composition Rules}

\begin{table}[htbp]
\centering
\small
\rowcolors{1}{white}{white}
\begin{tabularx}{\textwidth}{>{\centering\arraybackslash}X >{\centering\arraybackslash}X >{\centering\arraybackslash}X}
\toprule
\rowcolor{gray!25}\textbf{If you have...} & \textbf{Then you must also have...} & \textbf{Because} \\
\midrule
\rowcolors{1}{white}{gray!10}
$Ð_{\omega}$ & $Þ_{O}$ & Axiom C: self-referential dimensionality demands self-referential topology \\
$Ħ_{!}$ & $Ç^{\text{Ù}}$ or $Ç^{\lambda}$ & Axiom A: infinite memory requires kinetic trapping \\
$\Omega_{2}$ & $Ħ_{A}$ or $Ħ_{!}$ & Axiom B: binary topological protection needs $\geq$ 2-step memory \\
$\Omega_5$ & $Ð_{\omega}$ & Non-Abelian braiding requires infinite-dimensional state space \\
$\Phi_{\}}$ and $$\odot$_ÿ$ together & $\rightarrow$ $O_{\infty}$ candidate & Frobenius gate + criticality gate \\
$$\odot$_3$ in a tensor product & $\rightarrow$ kills $$\odot$_ÿ$ & EP absorption: measurement destroys self-modeling \\
\bottomrule
\end{tabularx}
\end{table}

\subsubsection{The Bottleneck Primitives}

When climbing tiers, only \textbf{certain primitives} are the rate-limiting steps:

\begin{itemize}
    \item \textbf{$O_0 \to O_1$:} $\phi$ is the driver. You must reach $\phi_ÿ$-criticality.
    \item \textbf{$O_1 \to O_2$:} $D$ and $\Omega$ are the drivers. Go from point to surface, from trivial to $\mathbb{Z}_2$.
    \item \textbf{$O_2 \to O_2^\dagger$:} Only $D$. Stretch surface to infinite.
    \item \textbf{$O_2^\dagger \to O_{\infty}$:} Only $\Phi$. The Frobenius leap --- from asymmetric to $\mu \circ \delta = \text{id}$.
\end{itemize}

\begin{center}
\rule{0.5\textwidth}{0.4pt}
\end{center}

\subsection{Part VI --- Exercises for Internalization}

\subsubsection{Exercise 1: The Twelve-Beat Warmup}

Recite the twelve primitive sounds in order, one per breath:

\begin{quote}
DAH (dimension), THORN (topology), RHA (relation), PHI (parity),
EFF (fidelity), SEE (kinetics), GAM (scope), GUH (logic),
PHI-HAT (criticality), AYK (memory), SIG (composition), OH (winding).

\end{quote}

Do this three times. Each time, replace the keyword with one subtype:

\begin{itemize}
    \item DAH: dot $\rightarrow$ deck $\rightarrow$ deh-in-fee $\rightarrow$ dot-oh
    \item THORN: thin-six $\rightarrow$ thick $\rightarrow$ th-RÓH $\rightarrow$ th-EH-um $\rightarrow$ th-OH
    \item ...
\end{itemize}

\subsubsection{Exercise 2: The Tier-Climb Drill}

Pick any $O_0$ system (e.g., a classical gas: $\langle D_{\triangle}; T_\text{net}; R_\text{sup}; P_\text{asym}; F_{\ell}; K_\text{fast}; G_{\beth}; \Gamma_{\wedge}; \Phi_\text{sub}; H_0; 1:1; \Omega_Å \rangle$).
Climb step by step:

\begin{enumerate}
    \item Promote $\Phi_\text{sub} \to \phi_ÿ$ $\rightarrow$ now $O_1$
    \item Promote $D_{\triangle} \to D_C$ AND $\Omega_Å \to \Omega_{\mathbb{Z}_2} \to$ now $O_2$
    \item Promote $D_C \to D_{\infty} \to$ now $O_2^\dagger$
    \item Promote $P_\text{asym} \to P_{\pm}^{\text{sym}} \to$ now $O_{\infty}$
\end{enumerate}

\subsubsection{Exercise 3: The Axiom Check}

For any proposed tuple, verify:

\begin{itemize}
    \item Is $D_{\odot}$ paired with $T_{\odot}$? If one exists without the other, the tuple violates Axiom C.
    \item Is $H_{\infty}$ paired with $K_{\text{trap}}$? If not, Axiom A is violated.
    \item Is $\Omega_2$ paired with $H_2$ or $H_{\infty}$? If not, Axiom B is violated.
\end{itemize}

\subsubsection{Exercise 4: The Sound-Shape Game}

Close your eyes. Hear "th-RÓH." What shape appears? An X, a bowtie, a crossing. That's $T_{\bowtie}$.
Hear "phi-shurl." What happens? A mirror reflecting perfectly, no distortion. That's $P_{\pm}^{\text{sym}}$.
Hear "ayk-bang." Everything echoes forever. That's $H_{\infty}$.

The sound-shape association is the fastest path to fluency. The sounds are not arbitrary --- the phonemes encode the primitives' geometric character.

\begin{center}
\rule{0.5\textwidth}{0.4pt}
\end{center}

\subsection{Part VII --- The 4$\times$4$\times$... = 17.28 Million Types}

The total count of structural types is:


\[4 \times 5 \times 4 \times 5 \times 3 \times 5 \times 3 \times 4 \times 5 \times 4 \times 3 \times 4 = 17{,}280{,}000\]

Each cell in this product space is a \textbf{Frobenius address} --- an integer from 0 to 17,279,999. The crystal of types (\pilcrow{}64 of the grammar) maps every tuple bijectively to exactly one address.

Most of these types ($60\% = 10.37$M) are $O_0$ --- subcritical, no self-modeling, no topology. Only $8\% = 1.38$M reach $O_{\infty}$ --- Frobenius-special, self-modeling, critically balanced, topologically protected.

\textbf{The grammar is a funnel:} the easy types are at the bottom, the rare types at the top. Every promotion upward costs structure --- you pay in primitives, and the cost grows monotonically.

\begin{center}
\rule{0.5\textwidth}{0.4pt}
\end{center}

\emph{This guide records the complete mnemonic map of all 12 primitives, all 49 subtypes, their sounds, images, composition axioms, and tier structure as verified by crystal\_tier\_census (17.28M types, 400 tier cells) and crystal\_tier\_gap\_ladder (O₀$\rightarrow$O₁ distance 1.05, O₁$\rightarrow$O₂ distance 1.30, O₂$\rightarrow$O₂† distance 1.00, O₂†$\rightarrow$O$\infty$ distance 4.38).}


\end{document}
